\section{The complete original Taylor--Wiles limit system}
\label{sec:limit-system-v039}

This chapter proves the whole original assertion \(\mathrm{S4}\), including
uniform convergence of the local homomorphisms.  Its scope is
\[
 \forall\ell\text{ prime},\ L,\ G_0,\ T:\mathrm{TWGrp}(G_0),\qquad
 T.\mathrm{ProL}\ \Longrightarrow\ \mathrm{Nonempty}(\mathrm{Lim}(T)).
 \tag{LIM.1}
\]
Here \(L\), \(G_0\), and \(T\) are the original objects, not replacements
with fewer fields.  In particular the two finite sets of local groups,
their actual maps at every level, the distinguished involutions,
the characters \(\eta,\epsilon\), and the Frobenius and tame elements
are retained.  The hypothesis \(T.\mathrm{ProL}\) is precisely the
antecedent in \texttt{PatchedSupportCompletions.S4}, lines 15--16.
Constructing the matched quotient that supplies this antecedent is a
separate arithmetic descent step.

Write \(\Gamma_0=G_0.\Gamma\),
\(\Gamma_N=(T.\mathrm{GN}\,N).\Gamma\),
\(p_N=(T.h\,N).p\), and \(H=T.H\).
The maps \(p_N\) are continuous surjective homomorphisms, and
\(H\) is an open normal subgroup of \(\Gamma_0\).
Put \(H_N=p_N^{-1}H\).  A profinite group below means a compact Hausdorff
totally disconnected topological group.  The original compact,
totally disconnected group assumptions imply the Hausdorff condition,
as shown in LIM1.  Products of normal subgroups denote their join.

For an open subgroup \(A\) of a profinite group, write
\[
 X_\ell(A)=\operatorname{Hom}_{\mathrm{cont}}(A,\mathbf F_\ell).
\]
This is the source's \(\mathrm{HomL}\), written additively: a homomorphism
to the finite discrete group \(\mathbf F_\ell\) is continuous exactly
when its kernel is open.  The \emph{literal} source hypothesis
\(\mathrm{PhiCond}(Y,\ell)\) says that \(X_\ell(A)\) is finite for
\emph{every open subgroup} \(A\le Y\), not merely for \(Y\) or for
one selected local inertia group.  The proof below needs it at the
specific open normal subgroup \(f_0^{-1}H\).

All nineteen claims have complete human proofs at their printed scopes.
They use elementary finite group, compactness, and linear algebra
arguments; there is no new published-theorem input.  Source citations
locate the precise formal contracts.  They do not replace the proofs.

% BEGIN LIM1
\paragraph{LIM1. Profinite quotients and finite coordinates.}
\label{claim:LIM1}
The compact totally disconnected topological groups in the original
\(\mathrm{Loc}\) and \(\mathrm{Glob}\) have a basis at the identity of
open normal subgroups.  Every such subgroup has finite index.  If
\(K\triangleleft P\) is closed, then
\[
 K=\bigcap_{V\triangleleft_o P}KV,\qquad
 P/K\ \simeq\ \varprojlim_{V\triangleleft_o P} P/KV
 \tag{LIM.2}
\]
as topological groups.  The quotient is again profinite.  A continuous
surjection of profinite groups is a quotient map.

Here are the needed topological details.  The closure of a singleton is
connected, so total disconnectedness makes every singleton closed.
A topological group with this property is Hausdorff: for \(g\ne1\),
choose an identity neighborhood omitting \(g\), and then a symmetric
identity neighborhood whose square lies in it; the corresponding
translates separate two distinct points.

In a compact Hausdorff space the intersection of all clopen
neighborhoods of a point is connected.  To see this, if that intersection
were the disjoint union of two nonempty compact sets, separate them by
disjoint open sets.  Compactness of the complement of those open sets
gives a finite intersection of clopen neighborhoods of the point lying
in their union.  Its intersection with the open set containing the point
is then clopen and excludes the other compact set, a contradiction.
Thus in a totally disconnected compact Hausdorff space the intersection
is the singleton.  A second finite-intersection argument gives a clopen
neighborhood inside every prescribed neighborhood.

Choose a clopen identity neighborhood \(C\) inside a given identity
neighborhood.  Continuity of multiplication and finite covers of the
compact sets \(C\) and \(P\setminus C\) give an identity neighborhood
\(W\) such that every element of \(W\) preserves both sets.  Explicitly,
cover each compact set by finitely many neighborhoods \(U_i\) for which
\(W_iU_i\) lies in that same set, and intersect the \(W_i\).
The stabilizer \(E=\{g:gC=C\}\) is a subgroup containing an identity
neighborhood and hence is open.  Since \(1\in C\), it is contained in
\(C\).  Open cosets cover the compact group, so \(E\) has finite index;
its normal core is a finite intersection of conjugates and is open.
This proves the asserted basis and index statement.

If \(x\notin K\), closedness of \(K\) gives an open normal \(V\)
with \(xV\cap K=\varnothing\), and therefore \(x\notin KV\).
This proves the first equality in (LIM.2).  A coherent family of cosets
in the finite quotients \(P/KV\) has the finite-intersection property:
for finitely many \(V\)'s use their intersection.  Compactness gives an
element of all the cosets, unique modulo \(K\).  Hence the indicated map
is a continuous bijection from the compact quotient onto the inverse
limit, which is Hausdorff.  A continuous bijection from a compact space
to a Hausdorff space is a homeomorphism because it sends closed sets to
compact, hence closed, sets.  This also proves that a continuous
surjection of profinite groups is a quotient map.  In particular,
compatible homomorphisms into all these finite quotients define a
unique continuous homomorphism into \(P/K\).
The source uses these facts in \texttt{QuotProf.totallyDisconnected\_quot}
and in the finite-quotient constructions of \texttt{LimKer} and
\texttt{InLim}.
% END LIM1

% BEGIN LIM2
\paragraph{LIM2. The relative pro-\(\ell\) residual, with its exact ambient meaning.}
\label{claim:LIM2}
Let \(H\triangleleft_o\Gamma\), with \(\Gamma\) profinite, and define
exactly as the source does
\[
 P_\ell^\Gamma(H)=
 \bigcap_{\substack{U\triangleleft_o\Gamma\\
       \forall h\in H\ \exists e\ge0:\ h^{\ell^e}\in U}} U.
 \tag{LIM.3}
\]
It is a closed normal subgroup contained in \(H\).  It equals, as a
subgroup of \(\Gamma\), the intrinsic residual
\[
 \bigcap_{\substack{A\triangleleft_o H\\H/A\text{ a finite }\ell\text{-group}}}A.
 \tag{LIM.4}
\]
The quotient \(H/P_\ell^\Gamma(H)\) is pro-\(\ell\).  Moreover for every
\(V\triangleleft_o\Gamma\),
\[
 P_\ell^\Gamma(H)\le V
 \quad\Longleftrightarrow\quad
 \text{the image of }H\text{ in }\Gamma/V\text{ is an }\ell\text{-group}.
 \tag{LIM.5}
\]

In a finite group, having all element orders powers of \(\ell\) is
equivalent to having order a power of \(\ell\).  The nontrivial direction
uses the following elementary proof of Cauchy's assertion.  If a prime
\(p\) divides the order of a finite group \(Q\), cyclic rotation acts on
the \(|Q|^{p-1}\) tuples \((a_1,\ldots,a_p)\) with product \(1\).
Rotation preserves that condition, also in a noncommutative group.
Nonconstant orbits have size \(p\).  The number of constant tuples,
equivalently solutions of \(a^p=1\), is therefore divisible by \(p\).
It includes the identity and hence includes a nonidentity solution.
Thus every prime divisor of \(|Q|\) occurs as an element order.
This establishes the equivalence and identifies the qualifying
subgroups in (LIM.3) with those for which the image of \(H\) is an
\(\ell\)-group.

The family in (LIM.3) contains \(H\) and is closed under finite
intersections: the image of \(H\) in the quotient by an intersection
embeds in the product of its finite \(\ell\)-group images.
Its intersection is consequently closed, normal, and contained in \(H\).
For any qualifying \(U\), \(H\cap U\) occurs in (LIM.4).
Conversely, if \(A\triangleleft_o H\) and \(H/A\) is an \(\ell\)-group,
take the core of \(A\) in \(\Gamma\).  Only the finitely many cosets of
\(H\) give distinct conjugates, because \(A\) is normal in \(H\).
The core \(U\) is open and normal in \(\Gamma\), is contained in \(A\),
and \(H/U\) embeds in a finite product of conjugates of \(H/A\).
It therefore qualifies in (LIM.3).  These two inclusions prove (LIM.4).

For the forward implication in (LIM.5), compactness gives finitely many
qualifying \(U_i\) whose intersection is contained in \(V\).  Indeed,
otherwise their closed intersections with the compact set
\(\Gamma\setminus V\) would have the finite-intersection property,
contradicting \(P_\ell^\Gamma(H)\le V\).
Their intersection still qualifies, so the image in \(\Gamma/V\)
is an \(\ell\)-group.  The converse is the defining intersection.
Finally, an open normal subgroup of \(H/P_\ell^\Gamma(H)\) pulls back
to an open normal subgroup \(A\) of \(H\) containing the residual.
Its core \(V\) in \(\Gamma\) is open normal and still contains the
residual.  By (LIM.5), \(H/V\), and hence \(H/A\), is an \(\ell\)-group.
Every finite continuous quotient of \(H/P_\ell^\Gamma(H)\) is thus an
\(\ell\)-group, as asserted.
This is the relative interpretation of \texttt{PL}, not a replacement
by a residual formed in an unspecified ambient group.
% END LIM2

% BEGIN LIM3
\paragraph{LIM3. Functoriality and surjective images of the relative residual.}
\label{claim:LIM3}
Let \(f:\Gamma\to\Gamma'\) be a continuous homomorphism of profinite
groups, let \(H\triangleleft_o\Gamma\), and let any subgroup
\(H'\le\Gamma'\) satisfy \(f(H)\le H'\).  Then
\[
 f(P_\ell^\Gamma(H))\le P_\ell^{\Gamma'}(H').
 \tag{LIM.6}
\]
If, further, \(H'\) is open normal, \(f\) is surjective, and
\(H=f^{-1}H'\), this inclusion is equality.
In particular, if the target residual is trivial, \(f\) kills the
entire source residual, not only its characters.

To prove (LIM.6), pull back any qualifying open normal subgroup \(U'\)
from (LIM.3) for \(H'\).  Its preimage is open normal and qualifies
for \(H\), so every element of the source residual maps into \(U'\).
Intersect over \(U'\).  In the surjective case the map
\(H\to H'\) is onto.  The image \(f(P_\ell^\Gamma(H))\) is closed
because it is a continuous image of a compact subgroup, and is normal
in \(\Gamma'\).  The quotient of \(H'\) by that image is a continuous
quotient of the pro-\(\ell\) group \(H/P_\ell^\Gamma(H)\) from LIM2.
It is pro-\(\ell\): any finite continuous quotient pulls back to a
finite continuous quotient of the latter group.
Every element of the intrinsic residual of \(H'\) is trivial in all
these finite quotients, and hence in the quotient itself by LIM1.
Thus \(P_\ell^{\Gamma'}(H')\le f(P_\ell^\Gamma(H))\), proving equality.
This standalone lemma supplies the generic scope of
\texttt{PlMap.pl\_map}; it does not construct an arithmetic
\(\mathrm{TWQuot}\) or transport any determinant or Hecke ideal.
% END LIM3

% BEGIN LIM4
\paragraph{LIM4. Finite \(\ell\)-groups are detected by linear characters.}
\label{claim:LIM4}
If \(Q\) is a finite \(\ell\)-group and \(M<Q\) is a proper subgroup,
there is a nonzero homomorphism \(Q\to\mathbf F_\ell\) vanishing on
\(M\).

Choose a maximal proper subgroup \(J\) containing \(M\).  Let \(J\)
act by left multiplication on \(Q/J\).  Every orbit has size a power
of \(\ell\), while \(|Q/J|\) is divisible by \(\ell\).
Thus the number of fixed cosets is divisible by \(\ell\).
There is at least one, the identity coset, so there are at least
\(\ell\) fixed cosets.  A coset \(gJ\) is fixed exactly when
\(g^{-1}Jg=J\); consequently the fixed cosets form \(N_Q(J)/J\).
The normalizer strictly contains \(J\).  Maximality forces
\(N_Q(J)=Q\), so \(J\) is normal.
The nontrivial quotient \(Q/J\) has no nontrivial proper subgroup.
It is cyclic, since any nonidentity element generates it, and its
order is prime, since a cyclic group of composite order has a
nontrivial proper subgroup.  Its order is therefore \(\ell\).
An isomorphism \(Q/J\simeq(\mathbf F_\ell,+)\), composed with the
quotient map, proves the assertion.  This proves the finite character
criterion used in \texttt{BurnsideHom.burnsideHom} without importing
an additional generator theorem.
% END LIM4

% BEGIN LIM5
\paragraph{LIM5. The exact finite-character generator bound.}
\label{claim:LIM5}
Suppose \(A\) is pro-\(\ell\) and \(X_\ell(A)\) is finite.  If
\(|X_\ell(A)|\le b\), then \(A\) has a list of \(b\) topological
generators, allowing identity entries.  Hence an open extension of
\(A\) in a profinite group is topologically finitely generated.

Pointwise addition and scalar multiplication make \(X_\ell(A)\) a
finite-dimensional \(\mathbf F_\ell\)-vector space: finite
intersections of open kernels show that sums are still continuous.
Write its dimension as \(d\).  Evaluation at group elements spans
its dual, because a character annihilated by all evaluations is zero.
Choose \(a_1,\ldots,a_d\) whose evaluations are a basis.
Let \(B\) be the closed subgroup they generate.
If \(B\ne A\), choose \(x\notin B\).  By compactness of \(B\)
and the identity basis from LIM1, some \(V\triangleleft_o A\)
satisfies \(x\notin BV\).  The image of \(B\) in the finite
\(\ell\)-group \(A/V\) is proper.  LIM4 supplies a nonzero character
of \(A/V\) vanishing on that image.  Its pullback is a nonzero member
of \(X_\ell(A)\) whose evaluations at all \(a_i\) vanish, a
contradiction.  Thus the \(a_i\) topologically generate \(A\).

Since \(d\le\ell^d=|X_\ell(A)|\le b\), append \(b-d\) identities.
The case \(d=0\) is covered: the argument forces \(A=1\).
If \(A\) is open in a profinite group \(B_0\), append finitely many
coset representatives of \(B_0/A\); the resulting finite list
topologically generates \(B_0\).  When the extension is not normal,
the same assertion means representatives of its finitely many left
cosets, and the generation argument is unchanged.
Applied to \(H_N\), this gives exactly \(T.\mathrm{hb}\) entries
from the literal \(\mathrm{finN}\) bound; no logarithmic or stronger
bound is being assumed.  This is the scope of \texttt{PlGen.pl\_gen}.
% END LIM5

% BEGIN LIM6
\paragraph{LIM6. A single finitely generated quotient for every local map.}
\label{claim:LIM6}
Let \(Y\) be a profinite group with the original
\(\mathrm{PhiCond}(Y,\ell)\).  Let \(f_0:Y\to\Gamma_0\) be
continuous and put
\[
 A=f_0^{-1}H,\qquad P_Y=P_\ell^Y(A).
 \tag{LIM.7}
\]
Then \(Y/P_Y\) is topologically finitely generated.
If continuous maps \(f_N:Y\to\Gamma_N\) satisfy
\(p_Nf_N=f_0\), then every \(f_N\) kills this \emph{same} \(P_Y\).

The subgroup \(A\) is open normal.  LIM2 shows that \(A/P_Y\)
is pro-\(\ell\).  Pullback along the surjection \(A\to A/P_Y\)
injects its continuous characters into \(X_\ell(A)\), which is
finite by the actual \(\mathrm{PhiCond}\) hypothesis at \(A\).
LIM5 makes \(A/P_Y\), and then its finite-index extension
\(Y/P_Y\), topologically finitely generated.
For each \(N\), the equality \(p_Nf_N=f_0\) gives
\(f_N(A)\le H_N\).  By \(T.\mathrm{ProL}\),
\(P_\ell^{\Gamma_N}(H_N)=1\).
LIM3 therefore gives \(f_N(P_Y)=1\) for all \(N\).
The quotient \(Y/P_Y\) is constructed once from \(f_0\) and \(H\);
it does not depend on the level or on a finite target used later.
This proves the exact local finiteness and factorization supplied by
\texttt{Loc\_fin.loc\_fin}.
% END LIM6

% BEGIN LIM7
\paragraph{LIM7. A free profinite group on a prescribed finite list.}
\label{claim:LIM7}
For every finite number \(r\) there is a profinite group \(F_r^{\mathrm{pf}}\)
with \(r\) distinguished topological generators, and every assignment
of those generators to elements of a profinite group \(Z\) extends
uniquely to a continuous homomorphism \(F_r^{\mathrm{pf}}\to Z\).

Take the abstract free group \(F_r\), formed by reduced words, and
map it diagonally to the product of all its finite quotients
\(F_r/N\), where \(N\) ranges over finite-index normal subgroups.
Let \(F_r^{\mathrm{pf}}\) be the closure of the image.
The product of finite discrete spaces is compact: extend any
family of closed sets with the finite-intersection property to an
ultrafilter, select its unique value at each finite coordinate,
and note that every basic neighborhood of the resulting point
belongs to the ultrafilter.  The point lies in all the original
closed sets.  This uses the same maximal-filter principle spelled
out in LIM9.  The product is Hausdorff, and its finite coordinate
projections separate points; hence the closed subgroup is profinite.
The images of the free generators topologically generate it by
construction.  An assigned list in \(Z\) defines an abstract word
homomorphism.  Modulo each \(V\triangleleft_o Z\), this homomorphism
has finite image, so its kernel has finite index in \(F_r\).
It consequently extends to a continuous map from
\(F_r^{\mathrm{pf}}\) to \(Z/V\), through the corresponding coordinate
finite quotient.  These maps are compatible: they agree on the dense
word subgroup and the finite targets are Hausdorff.
LIM1 assembles them into a continuous homomorphism to \(Z\).
Two such homomorphisms agree on the dense word subgroup and hence
everywhere.  For \(r=0\), the construction is the trivial group.
This supplies the universal property used in \texttt{FreeProf} and
\texttt{GenExt}, with its distinguished generators explicit.
% END LIM7

% BEGIN LIM8
\paragraph{LIM8. One finite profinite cover for all original levels.}
\label{claim:LIM8}
There is one topologically finitely generated profinite group \(\Pi\),
continuous surjections \(\pi_N:\Pi\to\Gamma_N\), and \(ct\in\Pi\)
such that
\[
 \pi_N(ct)=c_N,\qquad
 p_N\pi_N(x)\equiv p_M\pi_M(x)\pmod H
 \quad\text{for all }N,M,x.
 \tag{LIM.8}
\]

The quotient \(Q=\Gamma_0/H\) is finite; choose representatives
\(r_q\in\Gamma_0\) for all \(q\in Q\).
By LIM2 and \(T.\mathrm{ProL}\), each \(H_N\) is pro-\(\ell\).
By the actual \(T.\mathrm{finN}\) and LIM5, choose a list of
\(b=T.\mathrm{hb}\) topological generators of each \(H_N\).
Let \(\Pi=F_{1+|Q|+b}^{\mathrm{pf}}\), with its generators divided
into one distinguished generator \(ct\), a list indexed by \(Q\),
and a list of length \(b\).
At level \(N\), send \(ct\) to the original \(c_N\), send the
\(q\)-generator to a lift of \(r_q\) under the surjection \(p_N\),
and send the last list to the chosen generators of \(H_N\).
LIM7 gives a continuous homomorphism \(\pi_N\).

Its image is compact and therefore closed.  It contains the chosen
topological generators of \(H_N\), so contains \(H_N\), and its
image in \(\Gamma_N/H_N\simeq Q\) contains every coset.
It is therefore all of \(\Gamma_N\).
The maps \(p_N\pi_N\) have the same values modulo \(H\) on the
finite generators: respectively \(c_0\), \(r_q\), and \(1\).
Here \(p_N(c_N)=c_0\) is the original \(\mathrm{Glob.Hom.c}\) field.
Uniqueness on a dense generated subgroup gives their equality modulo
\(H\) on every \(x\in\Pi\), proving (LIM.8).
No relation \(ct^2=1\) has been imposed on \(\Pi\); its square
will lie in the limit kernel, which is exactly what is required.
This constructs the whole common cover from \texttt{SysCover.sys\_cover}.
% END LIM8

% BEGIN LIM9
\paragraph{LIM9. One ultrafilter and exact stabilization of the kernel joins.}
\label{claim:LIM9}
Fix an ultrafilter \(\mathcal U\) on \(\mathbf N\) extending the
cofinite filter.  Such an ultrafilter exists by the maximal-filter
argument: order proper filters containing the cofinite filter by
inclusion, take unions of chains, and apply the maximal principle;
maximality decides every subset or its complement.
It contains every tail, so in the source filter order
\(\mathcal U\le\mathrm{atTop}\).
Put \(K_N=\ker\pi_N\) and define exactly the source limit kernel
\[
 K=\{x\in\Pi:\ 
       \forall V\triangleleft_o\Pi,\quad
       x\in K_NV\text{ for }\mathcal U\text{-almost every }N\}.
 \tag{LIM.9}
\]
Then \(K\) is closed normal and, for every \(V\triangleleft_o\Pi\),
\[
 K_NV=KV\quad\text{for }\mathcal U\text{-almost every }N.
 \tag{LIM.10}
\]
Also
\[
 K\le V\quad\Longleftrightarrow\quad
 K_N\le V\text{ for }\mathcal U\text{-almost every }N.
 \tag{LIM.11}
\]

There are finitely many subgroups of \(\Pi/V\).  Thus the finite
partition according to \(K_NV\) has exactly one cell in \(\mathcal U\).
Write its subgroup as \(S_V\), so \(V\le S_V\) and
\(K_NV=S_V\) on that cell.  If \(W\le V\), intersect the two
cells and take one index in their nonempty intersection to obtain
\[
 S_W\le S_V,\qquad S_WV=S_V.
 \tag{LIM.12}
\]
For fixed \(x\), membership of \(x\) in \(K_NV\) is eventually
equivalent to membership in \(S_V\).  Consequently (LIM.9) is exactly
\(K=\bigcap_V S_V\), a closed normal subgroup.

We claim \(KV=S_V\).  One inclusion follows from \(K\le S_V\)
and \(V\le S_V\).  For the other, let \(x\in S_V\).
The closed sets \(S_W\cap xV\), indexed by open normal \(W\),
have the finite-intersection property.  For finitely many \(W_i\),
take \(W=V\cap\bigcap_iW_i\); by (LIM.12),
\(S_WV=S_V\), so \(S_W\cap xV\ne\varnothing\), and it lies
in all those sets.  Compactness gives \(k\in K\cap xV\).
Then \(x\in KV\), proving the claim and (LIM.10).
Finally, \(K\le V\) is equivalent to \(KV=V\), which by the
eventual equality is equivalent to \(K_NV=V\) eventually.
This is (LIM.11).  This argument proves the full
\texttt{LimKer.limKer\_sup} conclusion with its exact uniformity:
for each \(V\) the equality holds on one set in \(\mathcal U\).
No common set for infinitely many \(V\)'s has been asserted.
% END LIM9

% BEGIN LIM10
\paragraph{LIM10. Uniform convergence of the outgoing homomorphisms.}
\label{claim:LIM10}
Set \(F_N=p_N\pi_N:\Pi\to\Gamma_0\).
There is a continuous surjective homomorphism
\(F:\Pi\to\Gamma_0\), with \(K\le\ker F\), such that
\[
 \forall W\triangleleft_o\Gamma_0,\quad
 \mathcal U\text{-almost every }N:\quad
 \forall x\in\Pi,\quad F_N(x)F(x)^{-1}\in W.
 \tag{LIM.13}
\]

If a profinite group has \(r\) topological generators, continuous
homomorphisms from it to a finite group \(Q\) are determined by
the generator values and therefore number at most \(|Q|^r\).
Equality on the dense word subgroup implies equality everywhere,
which proves this assertion without any assumption that the
topological generators generate abstractly.
Apply it to \(\Pi\) and \(\Gamma_0/W\).
The maps \(F_N\bmod W\) take values in a finite set of
\emph{whole homomorphisms}.  The ultrafilter therefore chooses one
homomorphism \(f_W:\Pi\to\Gamma_0/W\) to which they are equal,
on one set in \(\mathcal U\), for all \(x\) simultaneously.
For \(W'\le W\), the two choices are compatible by intersecting
their eventual sets.  LIM1 gives a unique continuous homomorphism
\(F\) with these finite coordinates.  This proves (LIM.13).

Each \(F_N\) is onto, so each \(f_W\) is onto: choose one index
from its nonempty eventual set.  The image of \(F\) is compact
and closed.  If it omitted an element of \(\Gamma_0\), an open
normal coset separating that element from the image would
contradict surjectivity of one \(f_W\).  Thus \(F\) is onto.
For each \(W\), let \(V=\ker f_W\triangleleft_o\Pi\).
Eventually \(F_N\bmod W=f_W\), so \(K_N\le V\).
By (LIM.11), \(K\le V\).  Intersecting over \(W\) and using
finite-quotient separation proves \(K\le\ker F\).
This proves the original \texttt{LimOut.exists\_lim\_hom\_out}
endpoint; its convergence is uniform in all elements of the cover.
% END LIM10

% BEGIN LIM11
\paragraph{LIM11. The limit quotient and the exact all-level comparison modulo \(H\).}
\label{claim:LIM11}
Let \(G_* =\Pi/K\), let \(\pi_o:\Pi\to G_*\) be the quotient
map, and let \(h_*:G_*\to\Gamma_0\) be the homomorphism descended
from \(F\).  Both maps are continuous and surjective.  They satisfy
\[
 \begin{split}
 &h_*\pi_o=F,\qquad h_*(\pi_o(ct))=c_0,\\
 &p_N\pi_N(x)\,(h_*\pi_o(x))^{-1}\in H
       \qquad\text{for every }N,x.
 \end{split}
 \tag{LIM.14}
\]
The uniform statement (LIM.13) is also the original
\(p_{\mathrm{lim}}\) statement for these same maps.

LIM1 makes the quotient profinite, and LIM10 gives the well-defined
descent.  Its continuity follows from the quotient-map property;
its surjectivity follows from that of \(F\).
For every \(W\triangleleft_o\Gamma_0\), choose an index in the
eventual set of (LIM.13) and apply it to \(ct\).
Since \(F_N(ct)=p_N(c_N)=c_0\), the images of \(F(ct)\) and
\(c_0\) agree modulo every \(W\), proving their equality.
By (LIM.8), \(F_N\bmod H\) is already independent of \(N\).
It equals the ultrafilter's selected finite-coordinate homomorphism
\(f_H\), which is \(F\bmod H\).
Thus (LIM.14) holds for every \(N,x\), not merely eventually.
This all-level distinction is needed by the subsequent common
\(\eta\)-kernel and Selmer comparison.
% END LIM11

% BEGIN LIM12
\paragraph{LIM12. Uniform finite local homomorphisms.}
\label{claim:LIM12}
Keep \(Y,f_0,f_N,P_Y\) from LIM6 and the common cover and kernel
from LIM8--LIM11.  For every \(V\triangleleft_o\Pi\), put
\(S_V=KV\).  There is a continuous homomorphism
\[
 a_V:Y\longrightarrow\Pi/S_V
 \tag{LIM.15}
\]
with the following property.  On one set in \(\mathcal U\),
\(K_NV=S_V\), the finite quotient map \(q_V:\Pi\to\Pi/S_V\)
descends through \(\pi_N\) to a continuous homomorphism
\(q_{N,V}:\Gamma_N\to\Pi/S_V\), and
\[
 q_{N,V}f_N(y)=a_V(y)\qquad\text{for every }y\in Y.
 \tag{LIM.16}
\]
The maps \(a_V\) are compatible under refinement of \(V\).

On the set from (LIM.10), \(K_N\le S_V\).
Surjectivity of \(\pi_N\) gives the unique algebraic descent
\(q_{N,V}\).  It is continuous because \(\pi_N\) is a quotient
map and \(q_{N,V}\pi_N=q_V\) is continuous.
LIM6 says that \(f_N\) kills the same \(P_Y\) for every \(N\),
so \(q_{N,V}f_N\) factors continuously through \(Y/P_Y\).
That quotient is topologically finitely generated.
There are only finitely many continuous homomorphisms from it
to \(\Pi/S_V\), by the generator argument of LIM10.
For indices outside the kernel-stabilization set, use the trivial
homomorphism solely to make a sequence defined at every index.
The finite partition into these homomorphisms has a unique cell
in \(\mathcal U\); its homomorphism, pulled back to \(Y\), is
\(a_V\).  Intersecting that cell with the stabilization set gives
(LIM.16), simultaneously for all \(y\).

If \(W\le V\), then \(S_W\le S_V\).
Choose an index in the intersection of the two stabilization sets
and the two sets in (LIM.16).  The finite quotient maps at that
index commute, because their composites with the surjection
\(\pi_N\) are the original quotient maps of \(\Pi\).
Their composites with \(f_N\) therefore prove that \(a_W\)
maps to \(a_V\).  This is compatibility of the whole
homomorphisms.  It does not interchange an infinite intersection
of eventual sets with the quantifier over \(y\).
% END LIM12

% BEGIN LIM13
\paragraph{LIM13. Continuous local limits with the exact uniform \(\mathrm{Close}\) statement.}
\label{claim:LIM13}
There is a continuous homomorphism \(F_Y:Y\to G_*\) such that
\[
 \begin{gathered}
 h_*F_Y=f_0,\\
 \forall V\triangleleft_o\Pi,\quad
 \mathcal U\text{-almost every }N:\quad
 \forall y\in Y,\quad
 \mathrm{Close}(\pi_N,\pi_o,V)(f_N(y),F_Y(y)).
 \end{gathered}
 \tag{LIM.17}
\]
The source definition here is exactly
\[
 \mathrm{Close}(\pi_N,\pi_o,V)(a,b)
 \iff \exists x,x'\in\Pi:\
       \pi_N(x)=a,\quad\pi_o(x')=b,\quad x^{-1}x'\in V.
 \tag{LIM.18}
\]

By LIM1 and LIM9,
\(G_*\simeq\varprojlim_V\Pi/S_V\).
The compatible continuous homomorphisms in LIM12 therefore give
a unique continuous homomorphism \(F_Y\) with coordinates \(a_V\).
On the eventual set of LIM12, choose lifts \(x,x'\in\Pi\) of
\(f_N(y)\) and \(F_Y(y)\), respectively.
Equation (LIM.16) says that their images in \(\Pi/S_V\) agree.
Thus \(x^{-1}x'\in S_V=K_NV\).
Write \(x^{-1}x'=kv\) with \(k\in K_N\), \(v\in V\),
and replace \(x\) by \(xk\).
Its \(\pi_N\)-image is unchanged, while
\((xk)^{-1}x'=v\in V\).
This proves (LIM.18) for every \(y\) on the same eventual set.

It remains to prove the equality with the original map \(f_0\),
not merely an equality up to a new quotient.
Fix \(W\triangleleft_o\Gamma_0\) and put
\(V=\ker(F\bmod W)\).  Since \(K\le\ker F\), we have
\(S_V=KV=V\).  The map \(F\bmod W\) induces
\(\overline F_W:\Pi/V\to\Gamma_0/W\).
Choose one index in the intersection of the sets from LIM12 and
(LIM.13).  Then
\[
 \overline F_Wq_{N,V}=p_N\bmod W.
\]
Indeed, after composition with the surjection \(\pi_N\), this is
precisely \(F\bmod W=F_N\bmod W\).
Composing with \(f_N\) and using (LIM.16) gives
\(\overline F_Wa_V=f_0\bmod W\), since \(p_Nf_N=f_0\).
The left side is \(h_*F_Y\bmod W\).
Separation by all finite quotients of \(\Gamma_0\) proves
\(h_*F_Y=f_0\).
This is the full original \texttt{InLim.in\_lim} local endpoint,
together with the matching-map conclusion used by
\texttt{CloseMap.close\_map}.  It has been constructed through
finite homomorphism sets, not inferred from pointwise limits of
unrelated lifts.
% END LIM13

% BEGIN LIM14
\paragraph{LIM14. Both original local families converge on one common cover.}
\label{claim:LIM14}
For each \(v\in L.\mathrm{Vl}\), apply LIM13 with
\[
 Y=L.G_v(v),\quad
 f_N=\iota_{\ker\eta_N}\circ\phi_{N,v},\quad
 f_0=\iota_{\ker\eta_0}\circ\phi_{0,v}.
 \tag{LIM.19}
\]
The original \(T.\mathrm{phiV}\) gives \(\mathrm{PhiCond}(Y,\ell)\),
and the original \((T.h\,N).\phi\) gives \(p_Nf_N=f_0\).
All the continuity assumptions are the actual \(\mathrm{Glob}\)
fields, composed with the subgroup inclusion.
This constructs continuous \(F_v:L.G_v(v)\to G_*\), with
\(h_*F_v=f_0\), and the full uniform conclusion of (LIM.17).

Repeat the same construction for every \(u\in L.\mathrm{St}\),
using \(L.G_u(u)\), \(\phi^u_{N,u}\), \(T.\mathrm{phiU}\),
and \((T.h\,N).\phi^u\).  This gives continuous maps \(F_u\)
with their original composites to \(\Gamma_0\).
The same \(\Pi,\pi_N,K,\pi_o,h_*\), and \(\mathcal U\) are used
in both constructions.

For fixed \(V\triangleleft_o\Pi\), intersect the eventual sets
for the finitely many members of \(L.\mathrm{Vl}\), and separately
for the finitely many members of \(L.\mathrm{St}\).
This yields respectively
\[
 \begin{split}
 &\mathcal U\text{-almost every }N:\quad
   \forall v,y,\ \mathrm{Close}(\pi_N,\pi_o,V)
                 (\phi_{N,v}(y),F_v(y)),\\
 &\mathcal U\text{-almost every }N:\quad
   \forall u,y,\ \mathrm{Close}(\pi_N,\pi_o,V)
                 (\phi^u_{N,u}(y),F_u(y)).
 \end{split}
 \tag{LIM.20}
\]
An empty index set causes no exception: its assertion holds at every
level.  Each value of these maps lies in
\(\ker(\eta_0h_*)\), because its image under \(h_*\) is the
corresponding original value in \(\ker\eta_0\).
Codrestricting to this kernel is continuous in the subgroup topology.
This proves the exact \(\phi_{\mathrm{lim}}\) and
\(\phi^u_{\mathrm{lim}}\) fields, including the quantifier order
\(\forall V\), eventually \(N\), \(\forall\) place, \(\forall\)
local element.  None of the original \(I_v\) or \(J_u\) has been
changed; they remain the fields of the same \(L\).
% END LIM14

% BEGIN LIM15
\paragraph{LIM15. Limits of the actual Frobenius and tame elements.}
\label{claim:LIM15}
For every \(j\in\operatorname{Fin}(T.q)\), choose lifts
\(a_{N,j},b_{N,j}\in\Pi\) with
\[
 \pi_N(a_{N,j})=T.\mathrm{ph}_{N,j},\qquad
 \pi_N(b_{N,j})=T.\mathrm{ta}_{N,j}.
\]
There are elements \(a_j,b_j\in\Pi\) such that, for each
\(V\triangleleft_o\Pi\),
\[
 a_{N,j}^{-1}a_j\in V,\qquad b_{N,j}^{-1}b_j\in V
 \quad\text{for }\mathcal U\text{-almost every }N.
 \tag{LIM.21}
\]
To construct \(a_j\), use the unique ultrafilter-selected coset
of the finite quotient \(\Pi/V\) for the sequence \(a_{N,j}\).
These cosets are compatible by finite intersections of their
eventual sets, and LIM1 gives the element \(a_j\).
Do the same for \(b_j\).  This also proves (LIM.21).

Put \(\mathrm{ph}_j=\pi_o(a_j)\),
\(\mathrm{ta}_j=\pi_o(b_j)\).  The chosen lifts prove the exact
\(\mathrm{Close}\) assertions for these two elements.
Intersecting over the finite set of \(j\)'s proves both assertions
simultaneously for every \(j\), as the original \(\mathrm{Lim}\)
fields require.
Furthermore
\[
 h_*(\mathrm{ph}_j)=T.\mathrm{sg}_j,
 \qquad h_*(\mathrm{ta}_j)=1.
 \tag{LIM.22}
\]
For the first equality, fix \(W\triangleleft_o\Gamma_0\) and let
\(V=\ker(F\bmod W)\).
Intersect the set in (LIM.13), the set in (LIM.21), and the
eventual set supplied by the actual \(T.\mathrm{ph\_lim}\).
Its guard \(P_\ell^{\Gamma_0}(H)\le W\) holds because
\(T.\mathrm{ProL}\) makes that residual trivial.
The last set belongs to \(\mathcal U\) since it contains a tail.
At an index in their intersection, the successive congruences are
\[
 F(a_j)\equiv F(a_{N,j})
       \equiv p_N(T.\mathrm{ph}_{N,j})
       \equiv T.\mathrm{sg}_j\pmod W.
\]
All finite quotients separate points, giving the first equality.
For the second, the same argument uses the exact all-level field
\(p_N(T.\mathrm{ta}_{N,j})=1\).
Since \(T.\mathrm{sg}_j\in\ker\eta_0\), both limit elements lie
in \(\ker(\eta_0h_*)\).  They have not been replaced by arbitrary
elements with the same image in \(\Gamma_0\).
% END LIM15

% BEGIN LIM16
\paragraph{LIM16. The exact tame relation becomes commutation in the limit.}
\label{claim:LIM16}
For every \(j\), the two elements constructed in LIM15 commute:
\[
 \mathrm{ph}_j\mathrm{ta}_j
       =\mathrm{ta}_j\mathrm{ph}_j.
 \tag{LIM.23}
\]

Fix \(V\triangleleft_o\Pi\), let \(S_V=KV\), and let
\(M=|\Pi/S_V|\), a positive finite integer.
For \(\mathcal U\)-almost every \(N\), the equality
\(K_NV=S_V\) gives the continuous surjection
\(q_{N,V}:\Gamma_N\to\Pi/S_V\).
The original \(T.\mathrm{ta\_ker}\) says that
\(T.\mathrm{ta}_{N,j}\in H_N\).
The group \(H_N\) is pro-\(\ell\) by LIM2 and
\(T.\mathrm{ProL}\), so the image of this tame element in the
finite quotient has order \(\ell^e\) for some \(e\), with
\(\ell^e\le M\).  In particular \(e\le M\), including the
order-one case.  If \(N\ge M\), its order divides \(\ell^N\).

The original tame relation lies in
\(P_\ell^{\Gamma_N}(H_N)=1\), and therefore is the exact equality
\[
 T.\mathrm{ph}_{N,j}^{-1}\,T.\mathrm{ta}_{N,j}\,
 T.\mathrm{ph}_{N,j}
       =T.\mathrm{ta}_{N,j}^{\,T.\mathrm{Nq}_{N,j}}.
 \tag{LIM.24}
\]
The original congruence
\(T.\mathrm{Nq}_{N,j}\equiv1\pmod{\ell^N}\)
shows that its right side has the same image as
\(T.\mathrm{ta}_{N,j}\) in \(\Pi/S_V\) when \(N\ge M\).
Thus the images of the two actual level elements commute there.

Intersect the kernel-stabilization set, the tail \(N\ge M\),
and the two sets in (LIM.21).  At an index in this nonempty
intersection, the images of \(a_j,b_j\) in \(\Pi/S_V\)
are the images of those level elements, because \(V\le S_V\).
They therefore commute in \(\Pi/S_V\).
This holds for every \(V\), and \(\bigcap_V S_V=K\).
Their commutator is consequently in \(K\), proving (LIM.23)
in \(\Pi/K\).
The bound \(M\) is chosen after the finite quotient \(V\), as
permitted by the source; no threshold uniform over all finite
quotients is needed.  This argument does not claim that the
tame element itself is trivial in the limit.
It proves the exact \texttt{CommLim.comm\_lim} conclusion.
% END LIM16

% BEGIN LIM17
\paragraph{LIM17. Assemble every original \(\mathrm{Glob}\) and \(\mathrm{Hom}\) field.}
\label{claim:LIM17}
Define a \(\mathrm{Glob}(\ell,L)\) on \(G_*=\Pi/K\) by
\[
 \eta=\eta_0h_*,\qquad \epsilon=\epsilon_0h_*,\qquad
 c=\pi_o(ct),
 \tag{LIM.25}
\]
and use the kernel-valued maps from LIM14 for \(\phi\) and
\(\phi^u\).  This is an actual \(G:\mathrm{Glob}(\ell,L)\).
The same \(h_*\) gives an actual \(h:G.\mathrm{Hom}(G_0)\).

The group and topological fields, compactness and total
disconnectedness hold by LIM1 and LIM9.
The kernel of \(\eta\) is the preimage under the continuous
\(h_*\) of the original open \(\ker\eta_0\), so is open.
The homomorphism \(\epsilon\) is continuous because both factors
are continuous; its values are still the original
\(\mathbf Z_\ell^\times\), not a residual character.
For every \(N\),
\(\pi_N(ct^2)=c_N^2=1\), so \(ct^2\in K_N\) for all \(N\)
and therefore \(ct^2\in K\) by (LIM.9).
Thus \(c^2=1\).
LIM11 gives \(h_*(c)=c_0\), whence
\(\eta(c)=\eta_0(c_0)=-1\).
Both local families are continuous and kernel-valued by LIM14.
These are all the fields of the source \(\mathrm{Glob}\).

For \(h\), take \(p=h_*\).
Its continuity and surjectivity were proved in LIM11.
The \(\eta\) and \(\epsilon\) compatibility fields follow from
(LIM.25), the \(c\) field from LIM11, and the two local-map
compatibility fields from LIM13--LIM14.
These are all the fields of \(\mathrm{Glob.Hom}\), without any
change of local group or any new pointwise choice of global map.
Codrestriction gives the original \(h.p_0:G.G_0\to G_0.G_0\).
By LIM15, its values on the kernel-valued
\(\mathrm{ph}_j,\mathrm{ta}_j\) are exactly
\(T.\mathrm{sg}_j,1\).
The subgroups used by \(G.\mathrm{frame}\), including the same
\(L.I_v\) and the images of the same \(L.J_u\), are consequently
the literal source frame fields.
% END LIM17

% BEGIN LIM18
\paragraph{LIM18. Install the complete original \(\mathrm{Lim}(T)\) and prove S4.}
\label{claim:LIM18}
For every original \(T\) satisfying \(T.\mathrm{ProL}\), the
objects above form an element \(\Lambda:\mathrm{Lim}(T)\).
The full field verification is as follows; all entries refer to
one common \(\Pi,K,\mathcal U,\pi_N,\pi_o,G,h\).

\begin{longtable}{@{}p{0.34\linewidth}p{0.61\linewidth}@{}}
\toprule
Original fields & Constructed values and proof \\
\midrule
\endhead
\(\mathrm{Ph},\mathrm{grp},\mathrm{top},\mathrm{tgrp},
  \mathrm{cpt},\mathrm{tdisc},\mathrm{fg}\)
& \(\Pi\) from LIM8, with the finite list of free generators.
  The abstract subgroup it generates is dense, so its topological
  closure is the whole group, exactly the stated \(\mathrm{fg}\). \\
\(U,U_{\mathrm{le}}\)
& The single ultrafilter of LIM9, containing every tail. \\
\(\pi,\pi_{\mathrm{cont}},\pi_{\mathrm{surj}}\)
& The all-level continuous surjections of LIM8. \\
\(G,h\)
& The complete \(\mathrm{Glob}\) and \(\mathrm{Hom}\) of LIM17. \\
\(\pi_o,\pi_{o,\mathrm{cont}},\pi_{o,\mathrm{surj}}\)
& The same quotient \(\Pi\to\Pi/K\), by LIM11. \\
\(\mathrm{modH}\)
& The equality modulo \(H\) for \emph{every} \(N,x\) in LIM11. \\
\(ct,ct_N,ct_o\)
& The distinguished free generator, \(\pi_N(ct)=c_N\),
  and \(\pi_o(ct)=G.c\), by LIM8 and LIM17. \\
\(\mathrm{ker\_lim}\)
& LIM9: for each open normal \(V\), eventually
  \(\ker\pi_N\,V=K V=\ker\pi_o\,V\). \\
\(p_{\mathrm{lim}}\)
& LIM10--LIM11: for each open normal \(W\) of \(\Gamma_0\),
  eventually the comparison holds for \emph{all} \(x\in\Pi\). \\
\(\phi_{\mathrm{lim}},\phi^u_{\mathrm{lim}}\)
& LIM14: for each open normal \(V\), eventually for every
  original place and every element of its original local group. \\
\(\mathrm{ph},\mathrm{ta},\mathrm{ph\_lim},\mathrm{ta\_lim}\)
& The kernel-valued elements of LIM15; for each \(V\), their
  \(\mathrm{Close}\) assertions hold eventually for all \(j\). \\
\(\mathrm{comm}\)
& The exact equality in LIM16, also in the \(\eta\)-kernel subgroup. \\
\(p_{\mathrm{ph}},p_{\mathrm{ta}}\)
& The exact \(h.p_0\)-equalities of LIM15 and LIM17. \\
\bottomrule
\end{longtable}

The source's join notation is the product of the two normal
subgroups, so the displayed kernel identity is literal
\(\ker\pi_N\sqcup V=\ker\pi_o\sqcup V\).
The \(\mathrm{Close}\) relation is the original one in (LIM.18),
including its orientation \(x^{-1}x'\in V\).
Every object indexed by \(N\) still uses the given
\(T.\mathrm{GN}\), \(T.h\), \(T.\mathrm{ph}\),
\(T.\mathrm{ta}\), and \(T.\mathrm{Nq}\).
All intersections of eventual sets used to verify a single field
were finite.  The ultrafilter was selected once before the limit
maps and auxiliary elements, and no passage to a new subsequence
or a new ultrafilter occurred.
This proves (LIM.1), the whole original \(\mathrm{S4}\), for every
prime \(\ell\), including \(\ell=3\).
\(T.\mathrm{ker\_le}\), \(H_\epsilon\), \(\mathrm{fin0}\),
\(\mathrm{ph\_H}\), and \(\mathrm{sg\_mem}\) remain fields of
the supplied \(T\); the construction does not need to use them
to weaken or change the resulting object.
% END LIM18

% BEGIN LIM19
\paragraph{LIM19. Exact contribution to the selected Count-II support route.}
\label{claim:LIM19}
At the selected odd-prime Count-II scope, including \(\ell=3\),
suppose the original arithmetic data have been transported to one
matched \(T\) satisfying \(T.\mathrm{ProL}\), with their original
selected full-block point \(x\), \(a=n=4\), and chosen \(g=0\)
Kap \(k\).  Keep its original B1 bound with
\(D_0=4[L^+:\mathbf Q]+4T.q\) and its coupled chosen-Kap B2.
Then LIM18 supplies the actual previously missing
\(\Lambda:\mathrm{Lim}(T)\) input to CRW12--CRW15.
Those claims now give their selected K1, Laurent K1, K2 and K3
conclusions on this actual \(\Lambda.h\).

This is a substitution into their printed hypotheses, and its
same-object content is substantive.  The common cover, the all-level
\(\mathrm{modH}\) identity, the uniform outgoing comparison,
both uniform local-map comparisons, and the distinguished \(ct\)
are exactly the fields used in CRW12's finite-cocycle transport.
Its B1 constant is still chosen before the truncation \(m\);
for each \(m\), its permitted eventual set of levels belongs to
the \emph{same} \(\mathcal U\) because
\(\mathcal U\le\mathrm{atTop}\).
No original Selmer space is asserted to vanish, and no arbitrary
Kap B2 theorem is used.  The \(g=0\) simplification concerns
K2 and the coupled B2 only; the full \(4\times4\) first-order
coefficient space and the exact K1 bound remain.

The unconditional closure here is \(\mathrm{S4}\), with its
original \(T.\mathrm{ProL}\) antecedent.  Applying it to the
original arithmetic Taylor--Wiles system still requires the matched
\(\mathrm{TWQuot}\) descent, including the original determinant,
flags, old N4, point, B1/B2 and supplied \(\mathrm{TWAut}\).
LIM2--LIM3 provide generic group facts usable in that separate proof;
they do not assert that the original arithmetic \(T\) already has
trivial residual.  Nor does this chapter produce the matched
patched rings or modules, the S5 support dimension and fiber
identities, coefficient normalization, or an automorphic system.

Thus \(\mathrm{Raw}\) and \(\mathrm{COUNT\_II\_SUPPORT}\) remain
open until their other actual producers are supplied.  The original
ordinary/control/type data, four Hida endpoints, same-ideal
\(\mathrm{HasAut}\), complete CohP and LocCFT keep their original
scopes.  In the outer Count-II statement, the original complex
embedding is fixed first and one coefficient threshold serves both
rank-four Links and both Hida and P1 conclusions.  Each actual
coefficient ring retains its own original ideal \(J\); no ideal is
changed by constructing \(\Lambda\), and no common ideal across
different rings is asserted.  The same Count-II improvement serves
RH and unrestricted Dirichlet GRH; Case III alone has the previously
proved two-primary-target vacuity, while I, II and V remain substantive.
% END LIM19
